\begin{textsample}{NEG Dim 6 – pi – Score 1.00 – pi/pi\_inf\_e\_ef\_6.txt}  \label{ex:f6_neg_099}
2.
PRINCÍPIOS E CONCEPÇÕES Este Currículo está baseado em concepções pedagógicas que primam por uma educação de qualidade em que ensino e aprendizagem, por serem indissociáveis, têm estreita relação com o desenvolvimento pleno do estudante, conforme se discorre nos itens a seguir. 2.1 Educação : Conceito e abordagem A Lei de Diretrizes e Bases da Educação Nacional⁴ ( LDB, Lei nº 9.394/1996 ) concebe a educação como um processo formativo que ocorre em diferentes âmbitos de vivência dos sujeitos ( familiar, escolar, laboral, social e cultural ).
Essa concepção aponta para o entendimento da educação não como sinônimo de escolarização, mas de aprendizagens diversificadas e contínuas, que permeiam toda a vida dos indivíduos e dão respostas às suas diferentes questões.
A mesma lei, inspirada pela Constituição Federal de 1988 ( Art. 205 ), explicita que, especificamente, nas escolas, a educação tem como objetivo “ o pleno desenvolvimento do educando, seu preparo para o exercício da cidadania e sua qualificação para o trabalho ” ( BRASIL, 1996 ).
Em vigência, e pela orientação do Currículo do Piauí, essa concepção ampliada de educação e com foco no desenvolvimento pleno do estudante, pressupõe aprendizagens essenciais a serem mobilizadas e aplicadas nas diferentes esferas da vida.
Isso porque estamos imersos em um mundo, descrito por Delors et al. ( 1997 ), complexo e constantemente agitado, ao qual cabe à educação fornecer, de algum modo, os mapas e a bússola que permita aos estudantes navegar através dele.
Essa abordagem se alinha com a perspectiva do desenvolvimento de competências, indicada pela Base Nacional Comum Curricular ( BNCC/2017 ), que deve nortear as decisões pedagógicas nacionais no contexto do início do século XXI.
Ressalta -se que a BNCC define competência como “ mobilização de conhecimentos ( conceitos e procedimentos ), habilidades ( práticas cognitivas e socioemocionais ), atitudes e valores para resolver demandas complexas da vida cotidiana, do pleno exercício da cidadania e do mundo do trabalho ” ( BRASIL, 2017, p. 14 ).
Esse documento define dez competências gerais para a Educação Básica, que constituem os direitos de aprendizagem e desenvolvimento de todos os estudantes desse nível educacional ( vide quadro a seguir ) : ⁴ BRASIL.
Lei nº 9.394, de 20 de dezembro de 1996.
Estabelece as diretrizes e bases da educação nacional.
Disponível em :.
Acesso em : 28 mai. 2019.

% matched lemmas: 
\end{textsample}
